% !TeX encoding = UTF-8
\documentclass[variant=guia,cover=institutional,mode=screen]{ua-engineering-v6-article}
% !TeX encoding = UTF-8
\UASetUniversity{Universidad Austral}
\UASetFaculty{Facultad de Ingeniería}
\UASetDepartment{Departamento de Computación}
\UASetCampus{Pilar}
\UASetCareer{Ingeniería Informática}
\UASetCommission{Comisión A}
\UASetProfessor{Equipo docente}
\UASetSemester{Segundo cuatrimestre 2026}
\UASetCourse{Sistemas Distribuidos}
\UASetDate{2026}


\UASetClassNumber{03}
\UASetClassTitle{Consistencia y replicación}
\UASetDocKind{Guía de ejercicios}

\title{Clase 03 --- Guía}
\author{\UAProfessor}
\date{\UADate}

\begin{document}

\section{Indicaciones generales}
Resuelva cada ejercicio declarando supuestos, modelo de fallas, garantía y trade-off. Cuando corresponda, construya una ejecución verificable y vincule la respuesta con evidencia observable.

\section{Nivel 1: fundamentos y reconocimiento}
\begin{uaexercise}
Definir replicación y explicar sus motivaciones principales.
\end{uaexercise}

\begin{uaexercise}
Explicar por qué la divergencia es inevitable en un sistema asíncrono.
\end{uaexercise}

\begin{uaexercise}
Construir un ejemplo con dos réplicas que divergen tras escrituras concurrentes.
\end{uaexercise}

\begin{uaexercise}
Distinguir estado físico vs estado lógico en un sistema replicado.
\end{uaexercise}

\begin{uaexercise}
Explicar por qué “una sola fuente de verdad” es difícil de sostener en presencia de particiones.
\end{uaexercise}


\section{Nivel 2: ejecuciones y fallas}
\begin{uaexercise}
Definir linearizabilidad y dar un ejemplo que la viole.
\end{uaexercise}

\begin{uaexercise}
Definir consistencia causal con la relación happens-before.
\end{uaexercise}

\begin{uaexercise}
Explicar por qué dos timestamps físicos de nodos distintos no alcanzan para decidir causalidad entre eventos.
\end{uaexercise}

\begin{uaexercise}
Aplicar las reglas de relojes de Lamport a una secuencia simple de eventos con envío y recepción de mensajes.
\end{uaexercise}

\begin{uaexercise}
Dar un ejemplo permitido por causal pero no por linearizabilidad.
\end{uaexercise}


\section{Nivel 3: diseño y comparación}
\begin{uaexercise}
Definir consistencia eventual y sus límites.
\end{uaexercise}

\begin{uaexercise}
Comparar los tres modelos en términos de garantías observables.
\end{uaexercise}

\begin{uaexercise}
Describir leader-based replication y sus variantes sync/async.
\end{uaexercise}

\begin{uaexercise}
Explicar multi-leader y por qué introduce conflictos.
\end{uaexercise}

\begin{uaexercise}
Describir leaderless replication y el rol de $N,W,R$.
\end{uaexercise}


\section{Nivel 4: integración y defensa}
\begin{uaexercise}
Probar informalmente por qué $R+W>N$ implica intersección.
\end{uaexercise}

\begin{uaexercise}
Construir un ejemplo donde $R+W\le N$ rompe lecturas frescas.
\end{uaexercise}

\begin{uaexercise}
Definir conflicto concurrente y dar un ejemplo.
\end{uaexercise}

\begin{uaexercise}
Explicar LWW y sus riesgos.
\end{uaexercise}

\begin{uaexercise}
Explicar vector clocks y cómo detectan concurrencia.
\end{uaexercise}

\end{document}

% Material trasladado a la clase 4.
\iffalse
\section{Nivel 5: taller P--F--G--S--T sobre Asteria Online}

En los ejercicios siguientes use la notación del trabajo práctico:
\(P\)roblema, \(F\)alla, \(G\)arantía, \(S\)olución y \(T\)rade-off.
Mantenga la garantía formulada como un interés de negocio o de juego, sin
promesas absolutas ni números de implementación. En cada respuesta indique
qué observa el jugador y qué evidencia permitiría verificar la decisión.

\begin{uaexercise}
\textbf{Compra de oro con pérdida de señal.} Un jugador confirma una compra
de oro con dinero real y pierde la conexión antes de recibir la respuesta.
Construya un par \(P,G,F\): identifique el problema distribuido, formule una
garantía adecuada y describa al menos dos fallas de cliente relacionadas.
Explique por qué ``el jugador nunca paga dos veces'' está mal formulada como
garantía y cómo podría aparecer luego como propiedad de la solución.
\end{uaexercise}

\begin{uaexercise}
\textbf{Directorio de zonas.} Durante un particionamiento entre Brasil, US y
Zurich, dos datacenters reciben jugadores que intentan entrar a la misma
instancia. Proponga dos soluciones: una que priorice consistencia y otra que
priorice disponibilidad. Para cada una complete \(P,F,G,S,T\), describa una
ejecución posible y señale qué aspecto de CAP se prioriza.
\end{uaexercise}

\begin{uaexercise}
\textbf{Evento global.} El jefe tiene puntos de vida compartidos entre los
tres datacenters y recibe ataques concurrentes. Diseñe una solución para
evitar que el mismo ataque lógico se aplique dos veces cuando el cliente
reintenta. Luego describa una falla de seguridad que explote ese mecanismo y
la validación adicional necesaria del lado del servidor.
\end{uaexercise}

\begin{uaexercise}
\textbf{Recompensa única y mercado.} Dos jugadores obtienen o compran el
mismo objeto único durante una ventana de replicación atrasada. Compare
leader-based síncrono, leader-based asíncrono y leaderless con quórums.
Complete una tabla breve con la garantía protegida, el costo principal y la
condición bajo la cual cada alternativa deja de ser conveniente.
\end{uaexercise}

\begin{uaexercise}
\textbf{Caída de datacenter y migración.} Una mazmorra corre en Brasil y el
datacenter cae mientras el jefe está a punto de ser derrotado. Construya un
PFGST que diferencie el estado confirmado del estado en vuelo. Explique qué
parte de la garantía puede sostenerse, qué brecha queda durante la migración
y cómo se documenta esa brecha en el trade-off de la Entrega 3.
\end{uaexercise}

\begin{uaexercise}
\textbf{PFGST acumulativo.} Elija uno de los tres casos de uso de referencia
y produzca una tabla con al menos tres filas PFGST. Las filas deben incluir
una falla de cliente, una de infraestructura y una de seguridad; las
garantías deben permanecer iguales cuando sea razonable y las soluciones
deben evolucionar. Incluya una alternativa descartada en al menos una fila.
\end{uaexercise}

\begin{uaexercise}
\textbf{Garantía o promesa absoluta.} Clasifique y reformule estas frases
como garantías de negocio válidas: (a) ``el jugador siempre recibe la
recompensa'', (b) ``nunca se duplica el oro'', (c) ``el jugador no pierde una
compra de oro confirmada'', (d) ``el directorio responde en menos de dos
segundos''. Para una frase reformulada, indique cómo podría ser amenazada por
una falla de cliente, una de infraestructura y una de seguridad, sin cambiar
la garantía.
\end{uaexercise}


\fi
